\documentclass[11pt]{article}

% set this flag to 0 to remove comments
\def\comments{1}

% load the preamble
\usepackage{../preamble25}

\begin{document}

%%%%%%%%%%%%%%%%%%%%%%%%%%%%%%%%%%%%
%%%%%%%%%%%%  Title  %%%%%%%%%%%%%%%
%%%%%%%%%%%%%%%%%%%%%%%%%%%%%%%%%%%%

\HWtitle{2}{Wednesday, March 19, 2025}

%%%%%%%%%%%%%%%%%%%%%%%%%%%%%%%%%%%%
%%%%%%%%%%  Beginning  %%%%%%%%%%%%%
%%%%%%%%%%%%%%%%%%%%%%%%%%%%%%%%%%%%

\noindent
\textbf{Collaboration and Honesty Policy Reminder:}
Collaboration in the form of discussion is allowed. However, all forms
of cheating (copying parts of a classmate’s assignment, plagiarism
from papers or old posted solutions) are NOT allowed. A rough rule of
thumb: you should be able to walk away from a discussion of a homework
problem with no notes at all and write your solution on your own.
%
Finding answers to problems on the
Web or from other outside sources (these include anyone not enrolled
in the class) is forbidden. 

\begin{compactitem}
\item {\em You must write up each problem solution by yourself without
assistance, even if you collaborate with others to solve the
problem.}

\item You must identify your collaborators. If you did not
work with anyone, you should write ``Collaborators: none."

\item Asking and answering questions in every forum the class provides (on Piazza, in class, and in
  office hours) is encouraged!

\item Even though looking up answers is forbidden, using the web, generative AI, 
  or similar resources to find
  alternative explanations of concepts you need for the homework
  \emph{is} allowed, and encouraged. Asking Deepseek to solve a problem for you is not ok; asking it for an explanation of Chernoff bounds or for examples of encoding quardatic constraints in Gurobi is fine. You must \textbf{document your use of outside sources} and describe it at a high level in your solutions.
    \end{compactitem}

 \paragraph{Problems to be handed in}
\begin{enumerate}[leftmargin=\parindent, itemsep=3ex]
    \item \textbf{(In-class exercise from lecture 6)} Suppose you have a graph with a fixed vertex set $V$, and where
    each individual data point $x_i$ is an undirected edge $\{u,v\}\in V\times
    V$. For example, the nodes might represent locations, and an edge
    $\{u,v\}$ might represent the locations between which an individual
    travels most often.
  
    Consider the problem of finding a near-minimum cut in the graph. This
    is a partition of $V$ into two disjoint sets $A,B$ of nodes. The
    \emph{weight} of the cut is the number of edges that cross from $A$ to $B$
    (so $u\in A$ and $v\in B$ or vice versa).  The weight of a cut can be as large as the size of the data set $n$, and $n$ can be as large as $\Omega(|V|^2)$. 
    
    \begin{enumerate}
        \item Use the exponential algorithm (or report noisy max) to design an algorithm
        that returns a cut with expected weight $\text{\sf min-weight} +
        O(|V|/\eps)$.  \\ It's OK if your algorithm runs in time polynomial in  $2^{|V|}$.
        
        \item There can be multiple distinct minimum cuts in a graph.  However, one neat (and highly non-trivial to prove) fact is that if $w^* \geq 1$ is the number of edges in the minimum cut, the number of distinct cuts with weight $\leq c w^*$ is at most $O(|V|^{2c})$.  Using this fact, prove that the error of the exponential mechanism (or RNM) is actually much better than the bound in part (a). Namely,  it outputs a cut with expected weight
        $
        \text{\sf min-weight} +
        O(\log(|V|)/\eps)
        $.

        \item Give an $\eps$-differentially private algorithm $A$ that runs in time polynomial in $V$, with the following guarantee: if the minimum cut in the input graph $G$ has weight $w^* \geq \ln |V|$, then $A(G)$ returns a cut with weight $w^*+
        O(\log(|V|)/\eps)$ with probability at least $1-1/|V|$.
        % Mention in solutions that this is based on GLMRT '10.

    \end{enumerate}

    \item (\textbf{More on node-private graph analysis})
    Recall from Homework 1 that two graphs are \emph{node neighbors} if one can be obtained from the other by removing a node and all of its edges. Let $f(G)$ denote the number of triangles in an undirected graph $G$. You showed in the homework that, on the set of graphs with at most $n$ vertices, the global sensitivity of $f$ is $\binom {n-1} 2$.

    \begin{enumerate}
      \item Given a real-valued parameter $k>0$ and a graph $G$, consider the following linear program: 
      \begin{center}
      \begin{tabular}{|r l|}
          Variables: & $x_t$ for every triangle $t$ in $G$\\
          \hline
          Constraints: & For every triangle $t$, \qquad $0\leq x_t \leq 1$, and \\
          & for every node $u$, \qquad $\displaystyle \sum_{t: t \text{ contains }u} x_t \leq k$.\\
          \hline 
          Objective: & Maximize $\displaystyle\sum_t x_t$.
      \end{tabular}
      \end{center}
      Let $f_k(G)$ denote the value of this linear program (that is, the maximum possible value of the objective function). Show that 
      \begin{enumerate}
          \item Show that the value of this linear program equals the number of triangles in $G$ \emph{if and only if} every node $u$ is contained in at most $k$ triangles; and
          \item $f_k$ has global sensitivity $k$ (with no assumption on the neighboring inputs).
      \end{enumerate}

      So: one way to deal with the high sensitivity of $f$ is to instead release an approximation to $f_k$. In graphs where no node is involved in too many triangles, this will in fact approximate $f$.

    \item Let us now explore a different approach to dealing with the high sensitivity of $f$. 
    
    Given a graph $G$ and an integer $\lambda >0$, let $m_\lambda(G)$ denote the smallest value achieved by $f(H)$ on all subgraphs of $G$ that are obtained 
    by removing up to $\lambda$ nodes (and their edges) from $G$ (where $f$ is the number of triangles). 

    One way to deal with the high sensitivity of $f$ is, instead of aiming for an absolute error guarantee, to aim to release a value somewhere between $m_\lambda(G)$ and $f(G)$, for a value $\lambda$ that isn't too big. This means, roughly, that we are returning the number of triangles in ``almost all'' of $G$.



    \begin{enumerate}
      \item Show that increasing $\lambda$ decreases $m_\lambda(G)$. That is, $m_{\lambda+1}(G) \leq m_\lambda(G)$.
      \item Fix a parameter $\eps>0$ and a positive integer $\tau$. Consider the sequence
      $$\vec m(G) = (m_\tau(G), m_{\tau-1}(G), \dots , m_1(G), m_0(G))$$
      (where the last term $m_0(G)$ equals $f(G)$). Let $G'$ be a neighboring graph obtained by adding a new node along with an arbitrary set of edges to $G$. Show that the sequences $\vec m (G)$ and $\vec m (G')$ are interleaved, that is: 
      $$m_\tau(G) \leq m_\tau(G')\leq m_{\tau-1}(G) \leq  m_{\tau-1}(G') \leq \cdots \leq m_1(G) \leq m_1(G') \leq m_0(G) \leq m_0(G') \, .$$
      \item Suppose  we run the exponential mechanism to select a integer that is roughly a median of $\vec m(G)$ where $G$ is the input graph. More precisely, suppose we known a public upper bound $n$ on the size of $G$. Since $f$ takes values in $\cY = \{0,1,2,...,\binom n 3\}$, we use $\cY$ as our set of pissble outputs, and score function $\score(y;G) = |rank_{\vec m(G)(y)} - \frac \tau 2$. 
      
      Show that the score function has sensitivity 1 under insertion or deletion of nodes. Furthermore, show that if $\tau > 4\ln(n/\beta)/\eps$, then this mechanism will, with probability at least $1-\beta$ produce a value between $m_\tau(G)$ and $G$. 

    \end{enumerate}

    \end{enumerate}

    \item \textbf{(In class exercise from Lecture 6)} Show that the accuracy guarantees we showed for the
        exponential mechanism (and RNM) are basically tight in
        general. Specifically, consider the family of data sets
        $\set{\ds^{(1)},...,\ds^{(d)}}$ defined as follows: in
            $\ds^{(j)}$, one candidate $j$
        receives $\smax = n = \frac{\ln(d)}{2\eps}$  votes and all others
        receive 0 votes.
        \begin{enumerate}
        \item Show that on such inputs, the
        algorithm $A_{EM}$ will return a candidate other than $j$ (that is, a candidate 
        who received 0 votes)
        with constant probability, independent of $d$.
        \item \textbf{(*)}  Show that for \textbf{every}
        $\eps$-differentially private $A$
        algorithm, if we choose $J$ uniformly at random in $[d]$, then
        with constant probabililty $A(\ds^{(J)})$ will return a candidate
        other than $J$. [That is, $A$ will fail to find the winner for
        many datasets of this form.]
        \end{enumerate}

    \item (Programming problem, TBA)



        

  
  \end{enumerate}


\end{document}
